\section{Experiments and results}
\label{sec:results}

\subsection{Protocol}
\label{sec:protocol}

\paragraph{Two evidence sources, used for different things.}
This paper reports both simulated and real streams, and they answer different
questions. \Cref{sec:real-data} is the paper's empirical basis for what a real
backbone, on real surveillance footage, actually delivers: real per-frame
optical-flow scores from two public video-anomaly-detection corpora, real
train/test splits, real spliced anomalous events, and --- reported as a
finding rather than a footnote --- a real, severe autocorrelation structure
that leaves only a narrow band of usable operating points. It is also, by
construction, limited to the two datasets' native horizons (minutes, not
weeks) and to the one backbone we computed scores with. The simulated streams
that follow exist to answer the question the real corpora cannot: what
happens at the $10^5$--$10^6$-frame horizons, the wide range of episode counts
and intermittency levels, and the calibration/deployment shift severities a
municipal deployment must survive, but that no released corpus --- Mivia-IWDD
included --- provides months of continuously-labelled per-frame detector
output to test directly. Neither source substitutes for the other, and
\cref{sec:limitations} states plainly what each does and does not establish.

\paragraph{Simulated streams.}
SENTINEL-E consumes one scalar per frame, so the faithful way to stress the
guarantee layer at the timescales that matter --- weeks of video, $10^6$ frames
--- is to model that scalar. Our simulator is deliberately unfriendly in the ways
real footage is: a diurnal illumination cycle and a two-state weather chain shift
the null score distribution after calibration; scores follow an AR(1) with
Student-$t$ innovations, so they are autocorrelated and heavy-tailed; a
log-scale Ornstein--Uhlenbeck scene-activity process drives the backbone's
false-positive rate and is deliberately \emph{not} part of the operator's obvious
(clock, weather) taxonomy; dumping events are episodic and recurring, with
a configurable number of episodes per camera separated by exponential gaps, only
a fraction $\pi$ of the frames within an episode anomalous, and a ramp at each
end; and episodes propagate along the camera graph with a lag. Defaults put the null score near $0.24$ and
an anomalous frame near $0.75$, i.e.\ a per-frame signal weak enough that it
must accumulate over time --- the regime where sequential methods earn their
keep. \Cref{sec:limitations} states plainly what this does and does not
establish.

\paragraph{Baselines.}
A fixed score threshold with an optional $k$-of-$m$ persistence filter; a
per-frame threshold on the \emph{same} conformal p-values SENTINEL-E consumes,
which isolates optional stopping from calibration; CUSUM \citep{page1954};
Shiryaev--Roberts \citep{roberts1966}; a parametric e-detector using the same
changepoint mixture with a Gaussian likelihood-ratio bet instead of a conformal
one; a reimplementation of E-SHIFT \citep{khan2026eshift}; and, as the ablation that
isolates the paper's core, the same pipeline with the changepoint mixture in
place of the episodic one. That last baseline is not a straw man: by
\cref{prop:generalisation} it is a faithful instance of the PFA-controlled
non-partitioned e-detector of \citet{saha2026nonpartitioned} --- a mixture over
start times with weights summing to one, under a geometric prior and a fixed
betting family rather than an infimum over a composite pre-change class --- so
the episodes-versus-changepoints comparison in \cref{sec:episodes} is a
comparison against the current state of the art in anytime-valid change
detection, run through an identical conformal front end. Every baseline is
fitted on the same calibration set and evaluated on the same decorrelated
betting grid, so none is helped or penalised by the frame-rate choice.

Every method also \emph{restarts} after an alarm. This matters more than it
sounds. Under the single-shot convention of \cref{def:tuv} a detector stops at
its first alarm, so a method that false-alarms early is disabled for the rest of
the stream and records a missed detection for an event it never got to see. That
convention is right for stating the guarantee and wrong for measuring delay,
because an operator investigates an alert and resumes monitoring. Restarting
leaves the null metrics untouched --- the probability of at least one false
alarm, and the time to the first one, do not depend on what happens afterwards
--- and can only help a method that cries wolf. Under the single-shot convention
the classical baselines score miss rates near one; we report the restarting
numbers because they are the ones that give those baselines a fair hearing. Each is
swept over its own knob and plotted at its realised false-alarm rate.

\paragraph{Reproducibility.}
Seeds are fixed in each script; \texttt{experiments/run\_all.py} reproduces every
number, and every quantity quoted in this text is generated into a macro file
from the experiment output, so a stale number is not possible. All timings are
single-thread CPU on the machine that produced the statistical results.

\subsection{Real-video validation}
\label{sec:real-data}

Real video tests whether the guarantee holds per video on footage that the
simulator does not produce. Three corpora supply it: CUHK Avenue
\citep{lu2013avenue}, UCSD Ped2 \citep{mahadevan2010ucsd} and Mivia-IWDD
\citep{greco2026iwdd}, an illegal-dumping corpus. Each corpus is run with two
frozen backbones, and no backbone is fine-tuned.

\paragraph{Data and backbones.} Ped2 has \RealPedTwoNTrainVideos{} training
videos, all normal, and \RealPedTwoNTestVideos{} test videos from one fixed
camera. Avenue has \RealAvenueNTrainVideos{} training and
\RealAvenueNTestVideos{} test videos from one fixed camera. Mivia-IWDD has 200
training videos, all labelled negative, and 100 test videos: 50 negative, 25
static dumping and 25 dynamic dumping. Each Mivia video comes from a different
camera, so Mivia calibration pools cameras. The labels mark only the
dumping onset. Mivia is sampled at 10\,fps; the other two corpora at 25\,fps.

The flow backbone is Farneback optical-flow magnitude \citep{farneback2003},
averaged per frame at $64\times64$ resolution and median-filtered over three
frames within each video. The appearance backbone is PatchCore
\citep{roth2022total} with a frozen WideResNet-50-2 pretrained on ImageNet.
It uses layers 2 and 3, a 1\% coreset of training patches, and the maximum
over patches of the nearest-neighbour distance as the frame score. Its frame
AUCs are \RThreePedTwoFlowAuc{} (Ped2 flow), \RThreePedTwoPCAuc{} (Ped2 PatchCore),
\RThreeAvenueFlowAuc{} (Avenue flow), \RThreeAvenuePCAuc{} (Avenue PatchCore),
\RThreeMiviaFlowAuc{} (Mivia flow) and \RThreeMiviaPCAuc{} (Mivia PatchCore).
The Mivia flow score is below chance. PatchCore is a strong training-free appearance
backbone, but it is not the state of the art on these benchmarks. Its frame AUC is
0.828 on Ped2, 0.737 on Avenue (below flow's 0.783) and 0.588 on Mivia. The power
results below therefore bound what such frozen backbones deliver, not what the best
published video anomaly detectors would deliver. We did not add a state-of-the-art video anomaly detector. Such models are trained per dataset, and evaluating them without leakage would need its own training-only selection protocol.

\paragraph{Protocol.} The unit of replication is the video. Each test video is
monitored alone by a freshly started detector, and no frame is reused between
calibration and monitoring. Calibration uses training videos only. For Ped2
and Avenue it uses all training videos of the same camera. For Mivia it uses
the training videos of the other folds, pooled across cameras.

Operating points are selected on training videos only, before any test video
is read. SENTINEL-E's betting lag $L^\ast(\alpha)$ is the smallest lag on the
grid $\{1,2,3,5,7,10,15,20,30,50\}$ whose leave-one-video-out (LOVO)
per-video false-alarm proportion is at most $\alpha$. Ped2 and Avenue use
LOVO over their 16 training videos. Mivia uses 10 folds by video index, each
holding out 20 training videos. Each classical baseline is thinned at the same
$L^\ast(\alpha)$. Its knob is the first admissible knob in sweep order, from the
most aggressive. A knob is admissible when its LOVO proportion is at most
$\alpha$. If no knob is admissible, the most conservative knob is used. This
happens for \RThreeBaselineNoKnob{} of the \RThreeBaselineConfigs{} baseline
configurations at $\alpha=0.1$, and those results describe a conservative
setting, not a tuned one.

This replaces an earlier evaluation in which the betting lag was tuned on test
streams and normal frames were reused across spliced streams.

A null trial is a video with no anomaly, or the pre-onset part of an anomalous
video that lasts at least 2\,s. On Mivia the null trials are the 50 negative
test videos. A detection is an alarm inside an episode window, extended by a
one-second grace period. The detector restarts after each alarm, so a video
can contribute more than one alarm to the rate per hour. It counts once in the
per-video proportion. All proportions use exact Clopper--Pearson intervals.
False alarms per hour use an exact Poisson interval. The effective sample size
is stated with each result: the number of null videos, and the number of
episodes. Avenue has 74 episodes in 21 videos, so its detection intervals treat
episodes as independent and are optimistic.

\paragraph{Results.} \Cref{tab:r3_selection} gives the selected lags, and
\cref{tab:r3_test} the test results at $\alpha=0.1$. \Cref{fig:r3_per_video}
plots each method's test false alarms per hour against its detection rate.

\inputtable{tabr3_selection}

\inputtable{tabr3_test}

\begin{figure}[t]
  \centering
  \resultfig{fig_r3_per_video}{}
  \caption{Test results at $\alpha=0.1$, one panel per dataset and backbone.
  Each point is one method at its training-selected knob and $L^\ast(0.1)$.
  The horizontal axis is test false alarms per hour on a log scale, with zero
  drawn at $0.1$ per hour. The vertical bars are 95\% Clopper--Pearson
  intervals for the episode detection rate.}
  \label{fig:r3_per_video}
\end{figure}

The first finding concerns validity. On Ped2 and Avenue the false-alarm counts at
$\alpha=0.1$ are small. The null footage is short: 12.6\,s for Ped2 (three pre-onset
segments), 6.9\,min for Avenue (16 videos) and 26.9\,min for Mivia (50 negative
videos). At this length any outcome is consistent with the nominal level. Ped2 flow has \RThreePedTwoFlowFAk{}/\RThreePedTwoFlowFAn{} false
alarms, Ped2 PatchCore has \RThreePedTwoPCFAk{}/\RThreePedTwoPCFAn{}, Avenue
flow has \RThreeAvenueFlowFAk{}/\RThreeAvenueFlowFAn{} and Avenue PatchCore has
\RThreeAvenuePCFAk{}/\RThreeAvenuePCFAn{}. The intervals are wide. Ped2 flow at
$\alpha=0.2$ has 3 of 3 false alarms, with selected lag 1 (Online Resource~1, Table~S1),
so the Ped2 result depends on the level chosen. Mivia flow has
\RThreeMiviaFlowFAk{}/\RThreeMiviaFlowFAn{} false alarms, inside the level.
Mivia PatchCore has \RThreeMiviaPCFAk{}/\RThreeMiviaPCFAn{}, a proportion of
\RThreeMiviaPCFAp{} with interval $[\RThreeMiviaPCFALo{},\RThreeMiviaPCFAHi{}]$.
This exceeds $0.1$ at the point estimate, and the interval contains $0.1$. The
one-sided exact binomial probability of at least nine false alarms in 50 videos at
$p=0.1$ is $0.058$, so the excess is not significant at the 5\% level. The count also
moves with $\alpha$: it is 3 of 50 at $\alpha=0.05$ and 12 of 50 at $\alpha=0.2$
(full sweep in Online Resource~1, Sect.~S1). The excess has several candidate causes,
and 50 null videos cannot separate them. First, Mivia calibrates across cameras, but
the guarantee is per camera, and every test camera is unseen, so the per-camera
premise is not tested here. Second, the p-values are serially dependent (the
paragraph below). Third, the lag and the alarm boundary are chosen on training videos.
The selected lag is the smallest one whose LOVO proportion is at most $\alpha$. For
Mivia PatchCore, $L^\ast=7$ has 20 of 200 $=0.100$, exactly at $\alpha$, and selection at
a boundary can pick a lag whose true rate lies above $\alpha$.

Per-video rates understate what an operator sees, because the detector restarts after
each alarm. Per hour of null footage, SENTINEL-E raises 284.8 false alarms on Ped2
flow, from one alarm in 12.6\,s. On Avenue PatchCore it raises 86.6 per hour, and all
ten of its alarms come from one null video in 6.9\,min. On Mivia PatchCore it raises
100.5 per hour, with 45 alarms across nine videos in 26.9\,min. A per-video
proportion counts a video once, however many alarms it carries. The hourly figures
rest on very short null footage and are illustrations, not operating rates.

The second finding concerns power, which depends strongly on the backbone.
Ped2 flow detects \RThreePedTwoFlowDetk{} of \RThreePedTwoFlowDetn{} episodes
and Ped2 PatchCore detects \RThreePedTwoPCDetk{} of \RThreePedTwoPCDetn{}.
Avenue detects only \RThreeAvenueFlowDetk{} of \RThreeAvenueFlowDetn{} with
flow and \RThreeAvenuePCDetk{} of \RThreeAvenuePCDetn{} with PatchCore. Mivia
flow detects none of its \RThreeMiviaFlowDetn{} positive videos. Mivia
PatchCore detects \RThreeMiviaPCDetk{} of \RThreeMiviaPCDetn{}. Its median
delay is \RThreeMiviaPCDelayMed{}\,s after onset.

The third finding concerns the baselines, and it does not favour SENTINEL-E. On
Mivia PatchCore every baseline also exceeds $0.1$ on test, with 9 to 18 of 50
negative videos raising an alarm. On Avenue PatchCore, three of the six baselines
(CUSUM, Shiryaev--Roberts and the parametric e-detector) had no admissible knob at
$\alpha=0.1$ and used the most conservative setting. Across all 36 baseline
configurations at this level, 15 had none. The fixed threshold and E-SHIFT detect
more episodes than SENTINEL-E on Avenue PatchCore (15 and 29 of 74, against 9), with
2 and 4 false alarms in 16 videos. On Mivia PatchCore, CUSUM matches the 9 false
alarms in 50 videos, detects more episodes (11 of 50 against 10) and has a lower rate
per hour (53.6 against 100.5). Its knob is the most conservative in its sweep, so
this is not a tuned baseline. On Ped2 PatchCore the fixed threshold detects 8 of 12
episodes against 7 of 12, with no false alarm in three null videos, which is too
little evidence to rule out an excess. SENTINEL-E therefore shows no demonstrated
advantage over these baselines on real video. Its claimed advantage is the
horizon-free guarantee, which the short real videos cannot exercise.

\paragraph{Dependence.} The guarantee rests on the thinned p-values, so
\cref{tab:r3_dependence} measures their dependence. Two quantities are
reported. The pooled autocorrelation of the concatenated training scores
exceeds $0.05$ at every lag up to 200. That length reflects differences
between the mean scores of videos, not the time scale of the process, and
it is not used. The within-video autocorrelation is the relevant quantity. At
lag one it averages \RThreeWithinLagOneMin{} to \RThreeWithinLagOneMax{}. At
lag 25 it averages between essentially zero and \RThreeWithinLagTwentyFiveMax{}.
Lag 25 is 1\,s at 25\,fps and 2.5\,s at 10\,fps.

The thinned p-values remain positively correlated. At the selected lags the
lag-one correlation ranges from \RThreeDepCorrMin{} to \RThreeDepCorrMax{}. The
empirical CDF shows no excess of small values. The largest positive deviation
is \RThreeDepSupMax{} at the selected lags, and \RThreeDepSupOneMax{} at lag one.
The conditional-mean diagnostic contradicts conditional super-uniformity. If $p_t$
is conditionally super-uniform given the past, its mean given $p_{t-1}\le0.1$ is at
least $0.5$. The observed means lie between \RThreeDepCondMin{} and
\RThreeDepCondMax{}, on \RThreeDepCountMin{} to \RThreeDepCountMax{} pairs. A small
p-value is followed by another small one. The hypothesis of \cref{thm:episodic}
therefore fails on these thinned streams, and real-video validity can only be
established empirically. The observable consequence is the per-video false-alarm
rate. On Ped2 and Avenue it is consistent with $0.1$, but the null footage is short.
On Mivia PatchCore it exceeds $0.1$.

\inputtable{tabr3_dependence}

\paragraph{Action or residue.} The Mivia labels mark only the onset. They do
not say whether the detector responds to the dumping action or to the waste
left behind. Flow measures motion. Flow scores are higher on negative videos
than before onset in positive videos, so flow can at best respond to the
action, and the motion of other activity in the scene confuses it. Appearance
can in principle see deposited waste. Under PatchCore, static dumping is
detected in \RThreeMiviaPCStaticDetk{} of \RThreeMiviaPCStaticDetn{} videos and
dynamic dumping in \RThreeMiviaPCDynDetk{} of \RThreeMiviaPCDynDetn{}. These
counts do not differ at this sample size. The per-type split for all methods
is in Online Resource 1, Sect.~S1. A residue that persists after the action is
the permanent-change case, in which the episode never ends, $\eta=0$. The
episodic mixture contains that case. A dedicated residue model would change
the background model itself, and we leave it to future work.

\paragraph{Scope.} The real test videos last seconds to about a minute. The real data therefore bear only weakly on per-video false-alarm control, because
the null footage is short. The long-horizon
claim rests on the simulated streams in \cref{sec:validity}, which run for up
to a week of simulated time.

\begin{figure}[t]
  \centering
  \resultfig{fig_r0_qualitative}{}
  \caption{Real frames, real scores. Ped2 and Avenue rows: one normal and one
  anomalous frame, with the Farneback flow-magnitude field and its score, at
  the $64\times64$ resolution the pipeline uses. Mivia rows: the static and
  dynamic test videos, each with one pre-onset frame, one post-onset frame and
  the PatchCore score over time. Each video is the first positive test video of
  its type, in video-id order, with at least 3\,s before onset. It was not
  chosen on its score. The dashed line marks onset, and the coloured points mark
  the two frames. The frame panels are shown at display resolution, from the
  original releases. The Mivia release is not redistributed, so only downscaled
  stills are stored. The flow and score panels are computed from the pipeline
  inputs.}
  \label{fig:qualitative}
\end{figure}

The qualitative panels show how small the per-frame signal is. Ped2 moves from
\QualExamplePedTwoNormal{} to \QualExamplePedTwoAnomalous{}, and Avenue moves
from \QualExampleAvenueNormal{} to \QualExampleAvenueAnomalous{}. These are
the values the pipeline consumes, and they explain why the method accumulates
evidence over time.


\subsection{A single stream, worked through}
\label{sec:worked-example}

Every quantity in the sections that follow is an aggregate over hundreds of
simulated streams. Before reporting any of them, it is worth looking at one
stream in full, because the abstract objects in \cref{sec:method} --- a
p-value, a betting factor, accumulated wealth, an alarm --- are otherwise easy
to read only as symbols.

\begin{figure}[t]
  \centering
  \resultfig{fig6_worked_example}{}
  \caption{One simulated \WorkedT{}-second stream with \WorkedNEpisodes{}
  episodes (shaded), chosen for its episode count rather than for how
  favourably it behaves. \textbf{(a)} the raw backbone score, spiking during
  each episode but also, by construction, during ordinary null activity.
  \textbf{(b)} the conformal p-value at each betting instant (log scale);
  gaps are frames the decorrelation lag skips or the calibration abstains on.
  Notice that the score spikes in (a) do not translate into uniformly small
  p-values in (b) --- the score distribution the conformal layer calibrates
  against already absorbs much of the diurnal and weather-driven variation, so
  only the residual that context cannot explain shows up as evidence.
  \textbf{(c)} the episodic e-process's accumulated log-wealth. The dash-dot
  line is the labelled onset; the process is already drifting down slightly
  beforehand (an ordinary null fluctuation, licensed by the supermartingale
  property), climbs sharply once the first episode's evidence accumulates, and
  crosses $\log(1/\alpha)$ \WorkedDelaySec{}\,s after onset --- the red point.}
  \label{fig:worked-example}
\end{figure}

This particular stream is not cherry-picked for a fast detection: it is the
first seed tried at this configuration, and its three episodes are placed at
$\WorkedEp{}$. What it shows, concretely, is the mechanism the rest of the
paper measures in aggregate. The wealth process spends most of the stream
essentially flat --- a slow, licensed drift rather than a trend, since
$(W_t)$ is a supermartingale under the null and has no reason to grow --- and
climbs by more than four orders of magnitude in the roughly \WorkedClimbSec\,s
it takes the first episode's evidence to accumulate. After the alarm, further
episodes in this stream do not visibly move the trace: the wealth has already
crossed threshold and, in this single-alarm illustration, the detector is not
restarted, which is exactly the deployment convention (restart on alarm,
\cref{sec:protocol}) that the Monte Carlo experiments use and this one figure
does not. A reader who wants to see the restarted, multi-alarm behaviour on
this same stream, or to regenerate this figure on a different seed, can do so
with \texttt{experiments/exp11\_worked\_example.py}, which costs one stream
simulation and one pipeline run.

\subsection{False-alarm control is time-uniform}
\label{sec:validity}

\begin{figure}[t]
  \centering
  \resultfig{fig1_validity}{}
  \caption{Time-uniform false-alarm control, \ExpOneReps{} null streams per
  point in (a) and \ExpOneHorizonReps{} in (b). Open downward triangles mark
  operating points at which \emph{no} false alarm was observed; they are drawn
  at the 95\% Wilson upper confidence limit, which is the strongest statement
  the sample supports, and the dotted line connects them.
  \textbf{(a)} Realised against nominal level. SENTINEL-E stays below the
  diagonal everywhere; the per-frame p-value threshold, fed \emph{identical}
  conformal p-values, sits at the top of the plot, so the difference is
  attributable to the stopping rule alone. The middle curve is the same
  e-detector fed exactly uniform p-values, which separates the slack in Ville's
  inequality from the conservatism of the conformal layer.
  \textbf{(b)} Against the monitoring horizon at $\alpha=0.01$. The per-frame
  rules approach one; SENTINEL-E is flat.}
  \label{fig:validity}
\end{figure}

\inputtable{tab1_validity}
\inputtable{tab2_horizon}

\Cref{fig:validity} and \cref{tab:validity,tab:horizon} establish the basic
claim, and with \ExpOneReps{} streams per level the statement can be made
statistically rather than descriptively: at every level swept, the \emph{upper}
end of the $95\%$ interval on SENTINEL-E's realised false-alarm probability is
at most \ValidityWorstCiHi, and at the tightest level tested
($\alpha = \ValidityTightestAlpha$) it is \ValidityTightestCiHi{} --- below
nominal. The per-frame threshold on the same p-values reaches
\ValidityPfaPvaluePct\% at $\alpha=0.01$, with a mean of only
\ValidityPvalueArl{} frames --- \ValidityPvalueArlMin{} minutes of video --- to
its first false alarm. An operator running that rule across a hundred cameras
would field a false dispatch every few seconds. Extending the horizon from
\HorizonMinHours{} to \HorizonMaxHours{} hours drives the p-value threshold to
\HorizonPvaluePfa{} and the fixed score threshold from \HorizonFixedPfaShort{} to
\HorizonFixedPfa, while SENTINEL-E remains at \HorizonSentinelPfa.

SENTINEL-E is conservative, and it is worth saying where the conservatism comes
from rather than presenting it as free. Feeding the same e-detector exactly
uniform p-values yields a realised \ValidityPfaOracleLoose{} at a nominal
\ValidityAlphaLoose: most of the gap is slack in Ville's inequality and in the
changepoint prior, not a cost of the conformal layer. Ville's bound is not tight,
and a practitioner who wants a realised rate close to nominal should calibrate
the threshold empirically; what the theory provides, and what a fixed threshold
cannot, is that the bound never degrades with the horizon.

\subsection{Episodes versus changepoints}
\label{sec:episodes}

\begin{figure}[t]
  \centering
  \resultfig{fig5_episodes}{}
  \caption{What mixing over episodes buys, and where it does not, on the
  unpaired point estimates (\cref{tab:episodes}). With a single episode the
  episodic mixture is nominally behind rather than tied with the changepoint
  mixture --- \cref{tab:paired_episodes} resolves this as a real, if small,
  cost; the paired analysis below is what head-to-head claims in this section
  rest on.}
  \label{fig:episodes}
\end{figure}

\inputtable{tab11_episodes}
\inputtable{tab12_intermittency}
\inputtable{tab13_paired_episodes}

This is the experiment that isolates the paper's main contribution, and it
needed a second pass to actually isolate it. \Cref{tab:episodes,tab:intermittency}
sweep the two axes along which an episode differs from a changepoint, but their
two arms do \emph{not} share a betting grid: the episodic arm mixes over
\EpiGridKappa{} values of $\kappa$, the changepoint arm over $32$ log-spaced
values, so an unpaired point-estimate reading of these tables confounds the
episode-end axis with a grid-size effect that \cref{sec:ablation}'s own
``single bet aggressiveness'' row shows can be comparably large on its own.
\Cref{tab:paired_episodes} repairs this: both arms there run the identical
class over the identical \EpiGridKappa$\times$\EpiGridEta{} grid restricted to
$\eta=0$ against the full grid, differing only in whether $\eta$ ranges over
one value or four, and every seed is shared, so the delay difference is
reported as a paired bootstrap and the miss difference as an exact McNemar test
on discordant pairs rather than as two independent point estimates. Both arms
hold the false-alarm level throughout, as \cref{thm:episodic} and
\cref{prop:generalisation} require of both, so the comparison is entirely about
power.

The paired result is not monotone the way the unpaired tables suggested, and it
is more informative for not being so. At one episode the episodic mixture is
significantly \emph{behind}: mean censored-delay difference
\EpiPairedDiffOne\,s, $95\%$ CI \EpiPairedCIOne, bootstrap
$p=\EpiPairedPbootOne$, with the changepoint arm detecting
\EpiPairedDiscordantOne{} events the episodic arm misses and none of the reverse
(McNemar $p=\EpiPairedPMcnemarOne$). That is the price of the extra axis,
spending prior mass on episode-end hypotheses that never fire, and it is what
\cref{prop:generalisation} leads one to expect --- and it is a real, resolved
cost, not noise. At two and four episodes the sign reverses and is significant
in the same paired sense: at \EpiMidCount{} episodes the mean difference is
\EpiPairedDiffMid\,s in the episodic mixture's favour (CI \EpiPairedCIMid,
$p=\EpiPairedPbootMid$), with \EpiPairedDiscordantMid{} events caught only by
the episodic arm against none caught only by the changepoint arm
($p=\EpiPairedPMcnemarMid$). By \EpiMaxCount{} episodes the delay advantage is
smaller but still resolved (\EpiPairedDiffMany\,s, $p=\EpiPairedPbootMany$),
while the miss advantage --- \EpiPairedDiscordantMany{} discordant pairs out of
\ExpNineReps{} --- no longer clears significance on its own
($p=\EpiPairedPMcnemarMany$). None of this licenses the stronger, and wrong,
claim that the episodic mixture ``misses none at all'' in some meaningful
sense: \cref{sec:per-episode} shows that even where the \emph{stream}-level
miss rate is zero, the majority of individual episodes within that stream are
still missed, which is the mechanism this section's comparison is silent
about.

The crossover itself is the signature \cref{sec:episodic} predicts: it appears
exactly when there is more than one episode to accumulate over. Where the
mechanism claim needs qualifying is at the level of a single episode, which
\cref{sec:per-episode} checks directly.

\subsection{What the mechanism actually explains}
\label{sec:per-episode}

\inputtable{tab14_per_episode}

Every result so far scores a \emph{stream}: a run counts as detected once any
alarm follows the first onset, so a detector that sleeps through episodes $1$
through $k-1$ and only fires on episode $k$ is scored a clean detection. That
is not a hypothetical --- it is what happens here. \Cref{tab:per_episode} scores
every episode's own detection window independently (\ExpTenReps{} streams per
row) and the two columns diverge sharply: at \EpiMaxCount{} episodes the stream
miss rate is \PerEpiManyEpiStreamMiss{} for the episodic mixture, matching the
``misses none at all'' language earlier in this section, while the
\emph{per-episode} miss rate over the same streams is \PerEpiManyEpiMiss{} ---
the vast majority of individual episodes are never flagged inside their own
window. The stream-level statistic is not wrong, but it answers a narrower
question than ``does the detector catch dumping episodes'': it answers ``does
enough evidence eventually accumulate somewhere in the stream'', and with
\EpiMaxCount{} chances per stream, a per-episode catch rate as low as
$1 - \PerEpiManyEpiMiss{} \approx 0.13$ is compatible with a near-zero stream
miss rate.

The more surprising finding is that this per-episode miss rate is not where the
episodic mixture's advantage lives. At every episode count the episodic and
changepoint mixtures have statistically indistinguishable per-episode miss
rates --- \PerEpiOneEpiMiss{} against \PerEpiOneCpMiss{} at one episode,
\PerEpiMidEpiMiss{} against \PerEpiMidCpMiss{} at \EpiMidCount, and
\PerEpiManyEpiMiss{} against \PerEpiManyCpMiss{} at \EpiMaxCount. Neither
construction reliably resolves a single \EpiEventLenSec{}-second episode inside
its own narrow window; what differs between them, and what
\cref{tab:paired_episodes} establishes with a paired test, is how quickly
evidence that does not resolve within one episode's window compounds toward
the first stream-level alarm. The changepoint mixture carries a channel's worth
of burned wealth from a reverted bet into the next episode; the episodic
process's quiet state does not. That shows up as a shorter censored delay and,
at two and four episodes, as more streams resolved before the horizon --- both
already established in \cref{sec:episodes} --- but it is a claim about how fast
evidence accumulates \emph{across} a stream's episodes taken together, not a
claim that the episodic mixture independently catches more of them one at a
time. The contributions list and abstract should be read with that scope.

We report the single-episode row rather than omitting it. A strict
generalisation should cost something where its extra structure is absent, and
quantifying that cost is what makes the rest of the table trustworthy.

\Cref{tab:intermittency} varies the fraction $\pi$ of episode frames that show
the act, at \EpiBestCount{} episodes and with both mixtures using the same
undilated grid (\cref{sec:dilation}). Down to $\pi = 0.5$ the task is easy
enough that the two are indistinguishable and neither misses anything; the
separation appears only at $\pi = \EpiMinPi$, where the episodes become weak
enough that several are needed before the evidence crosses the threshold ---
which is, again, precisely when accumulating across episodes is what matters.
Below $\pi \approx 0.2$ neither variant detects reliably at any setting we
tried, and we take that as the honest edge of the operating envelope rather than
a tuning problem.

\subsection{Detection delay at a controlled false-alarm rate}
\label{sec:delay}

\begin{figure}[t]
  \centering
  \resultfig{fig2_delay_far}{}
  \caption{Detection delay against realised false-alarm rate,
  \ExpTwoReps{} null and \ExpTwoReps{} post-change streams per operating point.
  Lower and further left is better; the operating region a municipality can live
  with is the left-hand edge. Each method is swept over its own knob and plotted
  where it actually lands, not where it claims to.}
  \label{fig:delayfar}
\end{figure}

\inputtable{tab3_delay_far}

This is the comparison the paper exists to make, and to our knowledge no
existing waste-surveillance paper reports its analogue. \Cref{fig:delayfar} and
\cref{tab:delay_far} give delay against realised false-alarm rate for every
method. SENTINEL-E detects in \DelaySentinelSec\,s
($\pm$\,\DelaySentinelCi\,s) at a realised \DelaySentinelPfa, missing
\DelaySentinelMiss{} of events. \DelayNumInadmissible{} of the baselines never
reach a realised false-alarm probability of $0.05$ at any knob setting while
still detecting events (\DelayInadmissibleList), which is the substantive
finding: on streams with drift, heavy tails and autocorrelation, the classical
procedures cannot be tuned into the operating region at all, because the
assumptions their thresholds rest on do not hold.

\subsection{Robustness to calibration/deployment mismatch}
\label{sec:shift}

\begin{figure}[t]
  \centering
  \resultfig{fig3_shift}{}
  \caption{Shift in an \emph{unbinned} covariate (scene activity). Binning the
  context cannot help with a variable outside the taxonomy, so Mondrian
  calibration fails alongside the pooled baseline; conformalising the
  context-normalised residual holds throughout.}
  \label{fig:shift}
\end{figure}

\inputtable{tab4_coverage}
\inputtable{tab5_activity}

Two kinds of mismatch, both routine between the week a calibration set is
labelled and the month a camera is monitored, and the more interesting of the
two produces the clearest limitation in this paper.

\Cref{tab:activity} and \cref{fig:shift} shift an \emph{unbinned} covariate:
deployment scene activity is a multiple of the labelling window's. Pooled and
Mondrian calibration lose control immediately --- both reach
\ShiftPooledPfa{} and \ShiftMondrianPfa{} at the largest shift --- and the
Mondrian failure is instructive rather than embarrassing, since it is exactly
group-conditional on the variables it bins and activity is not one of them. An
operator cannot bin on every covariate that matters, which is the argument for
\cref{eq:residual}.

The residual construction does far better but it does not hold indefinitely,
and reporting where it breaks needs the same care as \cref{sec:ablation}: at
\ExpThreeReps{} replicates per shift level, a point estimate of PFA $0.000$ or
$0.02$ against nominal $0.01$ is not, by itself, evidence of validity or its
absence. Judged by the Wilson interval rather than the point estimate, shifts
of \ShiftResidualUnresolvedRatios{} are unresolved --- their intervals contain
nominal in both directions --- and only at $\times\ShiftResidualBreaks$ does
the lower end of the interval clear nominal, at a realised rate of
\ShiftBreakResidualPfa{} $\in[\ShiftBreakResidualPfaLo,\ShiftBreakResidualPfaHi]$.
That is the one shift level this design can certify as broken, not the
smallest shift at which it might be. We report this plainly because it is the
honest boundary of the method: normalising by a fitted $\hat g$ and
$\hat\sigma$ absorbs a shift in $P(c)$ only while the model stays adequate,
and a large enough move in an unmodelled direction breaks it, but pinning down
exactly how large would need many more replicates than run here.

This is where the likelihood-ratio weighting of \cref{sec:layer1} earns its
place, and where our own expectation was wrong. Weighting holds the realised
rate at \ShiftBreakWeightedPfa{} at every shift level tested, including the ones
where the plain residual construction fails --- but it does so by becoming
drastically conservative, with the miss rate rising from
\ShiftBreakResidualMiss{} to \ShiftBreakWeightedMiss{} at
$\times\ShiftResidualBreaks$. It buys validity with power, which is the correct
trade for a guarantee and a poor one for a detector.

The practical recommendation follows directly and is implemented in the library:
run the classifier two-sample test continuously, use the plain residual
construction while it reports no detectable shift, and enable weighting once it
does --- accepting that the camera has become a conservative detector and that
the operator should be told the calibration set no longer covers the conditions.
An abstention or miss rate that jumps is a signal to recalibrate, not a defect
to hide.

\Cref{tab:coverage} varies coverage instead. Validity survives all three
regimes, but power does not, and it should not: with calibration from a single
daytime window the miss rate reaches \CoverDaytimeMiss{} and the abstention rate
\CoverDaytimeAbstain\%. Degradation appears in the abstention and miss columns,
where an operator can see it, rather than silently in the alarm rate, where they
cannot. The operational implication is a requirement rather than a caveat: the
calibration set must span the conditions you intend to monitor, and the
abstention rate tells you whether it does.

\subsection{The camera network}
\label{sec:network}

\begin{figure}[t]
  \centering
  \resultfig{fig4_network}{}
  \caption{Fleet of \NetCameras{} cameras over \NetFleets{} held-out fleets with
  events and \NetFleets{} all-null fleets. The learned spatial prior does
  \emph{not} shorten delay --- it is worse than no graph at all, the negative
  result detailed below --- while e-BH holds the false-discovery rate under
  the dependence the graph induces regardless of which controller feeds it.}
  \label{fig:network}
\end{figure}

\inputtable{tab6_network}

\Cref{tab:network} separates the two questions, and they get different answers.

\paragraph{The guarantee survives learned control, as claimed.}
This is the part \cref{thm:predictable} is about, and it holds empirically as
well as formally. On all-null fleets every controller --- none, hand-designed,
and learned --- produces an alarm rate of \NetGnnNullPfa{} and an e-BH
false-discovery rate of \NetGnnFdr{} against a target of \NetFdrLevel. The
learned controller does this while actively \emph{hurting} power, which is the
sharpest available demonstration of the theorem: a badly-behaved learned
component degrades detection and leaves the false-alarm guarantee untouched,
because it can only enter through a predictable modulation.

\paragraph{The learned spatial prior does not earn its place.}
We expected the graph to shorten delay and it does not, on either delay
statistic: conditional delay is \NetHeuristicDelaySec\,s (heuristic) and
\NetGnnDelaySec\,s (learned) against \NetNoGraphDelaySec\,s with no graph, and
censored delay --- the comparable statistic per \cref{sec:problem}, since the
three controllers' miss rates differ --- is \NetHeuristicCensDelaySec\,s and
\NetGnnCensDelaySec\,s against \NetNoGraphCensDelaySec\,s. The hand-designed
prior is indistinguishable from no graph at all on either statistic, with power
\NetHeuristicPower{} against \NetNoGraphPower. The learned controller is worse
than both on every axis: delay, censored delay, and power \NetGnnPower.

We made two attempts to fix this before concluding it, and the second is worth
recording because it failed in an informative direction. The first was
initialisation: the controller's heads are set so that an untrained network
reproduces the ungoverned detector exactly, rather than starting at half stake
and a twenty-five-fold hazard, so training begins from the baseline and can only
improve on it in-sample. That is now the default and it is what the table
reports. The second was the objective. Maximising terminal log-wealth is
plainly the wrong target --- detection delay is the time the wealth takes to
reach $1/\alpha$, so a camera that ends the stream rich but climbed slowly is no
use --- and we replaced it with log-wealth credited over a
\NetCreditWindowSec\,s window opening at each camera's own onset, trained and
evaluated exactly as the terminal-objective controller was. The
better-motivated objective performed \emph{worse}: power fell from
\NetGnnPower{} to \NetGnnWindowedPower. We report the stronger of the two
configurations in \cref{tab:network} and include both there, and keep the
windowed objective in the library behind a flag so the comparison reproduces.

That a better-motivated objective makes things worse, and that the training gain
under it drifts downward rather than up, both point the same way: the gradient
signal is weak, not the optimiser mistuned. There is little for the controller
to learn.

Our reading is that the spatial signal in these streams is real but small
relative to what a camera learns from its own p-values: an episode propagates to
a neighbour with probability $0.45$ after a lag of order a thousand frames, by
which time the neighbour's own evidence usually dominates. A denser fleet,
stronger propagation, or a site layout where a single access road feeds many
cameras might change that, and we have not tested those. What we can say is that
on the configuration reported here the graph layer's value lies entirely in
e-BH, not in the spatial prior.

\paragraph{Fleet-level FDR control does earn its place.}
e-BH recovers the affected cameras at power \NetNoGraphPower{} with a realised
false-discovery rate of \NetNoGraphFdr{} against a target of \NetFdrLevel,
under dependence induced by shared weather, shared lighting and correlated
episodes. That is the part of Layer 3 we would keep.

\subsection{Computational cost}
\label{sec:compute}

\inputtable{tab7_compute}

\Cref{tab:compute} measures what an added-cost claim can be measured to mean:
a same-CPU ratio against a stand-in backbone we built to the scale of a real
edge detector, not a measurement on a specific published model or deployed
hardware, and one that excludes the cost of extracting the low-dimensional
context features Layer 1 consumes (time of day, a weather flag, a
background-subtraction motion measure) and the peak memory the guarantee layer
needs, neither of which this table reports. Within that scope: the reference
backbone --- a depthwise-separable stack with \CostBackboneParams\,M
parameters at $320\times320$, timed on the same CPU --- costs
\CostBackboneMs\,ms per frame. SENTINEL-E adds \CostSentinelUs\,\textmu s per
frame amortised, or \CostSentinelPct\% of the backbone, running at
\CostRealtimeFactor$\times$ real time on one thread. Fitting a camera's
calibration takes \CostFitSec\,s once. The reason the added cost is this small
is structural rather than incidental: the per-frame arithmetic is the $O(1)$
forward recursion of \cref{eq:forward}, and it runs once per decorrelation lag
rather than once per frame. Hardware that already runs the detector in real time,
and already extracts the context features, needs no additional
budget.

\subsection{Checking the assumptions}
\label{sec:diagnostics}

\inputtable{tab9_diagnostics}
\inputtable{tab10_pcoverage}

The guarantee rests on two premises that are assumptions rather than theorems.
The first, conditional independence, is checkable from the calibration set
alone, so an operator can run it on their own camera before trusting the alarm
rate; the second needs labelled null deployment data an operator does not
normally have, and we report it here as a check on the library rather than a
field-deployable one. \Cref{tab:diagnostics} reports conditional independence:
betting on every frame leaves a lag-one residual autocorrelation of
\DiagAcfRaw{} and a Ljung--Box p-value of \DiagLjungRaw, which rejects the
premise outright, while betting once per estimated lag gives \DiagAcfThin{}
and \DiagLjungThin{} --- markedly better, but \cref{tab:diagnostics}'s own
pass-rate row shows this reduces the violation rather than removing it: see
\cref{sec:limitations}. \Cref{tab:pcoverage} checks Layer 1 in isolation on
$191{,}890$ betting instants pooled from labelled null streams: the deployment
p-values are super-uniform at every level tested, with a smallest attainable
value of \DiagPmin{} --- the tail resolution that \cref{thm:beta} buys and that
the DKW inflation would forfeit --- but this is a validation of the library
against ground truth we constructed, not a check an operator can run on an
unlabelled camera.

\subsection{Ablation}
\label{sec:ablation}

\inputtable{tab8_ablation}

\Cref{tab:ablation} removes one decision at a time, and reports a Wilson 95\%
interval on the realised PFA rather than the point estimate alone: at
\AblReps{} replicates per row, far fewer than the \ExpOneReps{} used for the
validity sweep in \cref{sec:validity}, most rows reporting PFA $0.000$ have a
Wilson upper bound above nominal and are honestly unresolved, not certified
valid. \AblNumInvalid{} of \AblNumVariants{} variants are nonetheless \emph{clearly}
invalid --- their interval's lower end already exceeds $\alpha$ --- and the
pattern the design was built to show survives that correction: every row whose
invalidity is resolvable is a temporal or calibration one (betting every frame,
\AblEveryFramePfa; betting once every five frames rather than at the estimated
decorrelation lag; pooled conformal, \AblPooledPfa; Mondrian conformal;
omitting the calibration-conditional correction, \AblNoCcPfa), while no row
that only changes the betting family or the onset prior is ever resolved as
invalid. \AblNumUnresolved{} rows --- among them raw (unthinned) calibration,
the DKW inflation, and residual conformal with likelihood-ratio weighting ---
are unresolved rather than confirmed valid at this replication count; a point
estimate of $0.010$ for unthinned calibration in particular is not evidence
that skipping the thinning is safe, only that \AblReps{} replicates cannot
show it is not. Properly resolving these would need the replication count of
\cref{sec:validity} (\ExpOneReps{}), run there for the episodic detector alone;
see \cref{sec:limitations}.

Three rows deserve comment regardless. Replacing the episodic mixture with the
changepoint mixture costs power --- \AblChangepointMiss{} of events missed
against \AblFullMiss, with censored delay \AblChangepointDelay\,s against
\AblFullDelay\,s, on the \emph{same} \AblReps{} seeded streams as the
\EpiMaxCount{}-episode row of \cref{sec:episodes} rather than an independent
replication of it --- while both remain in the unresolved class on validity,
not confirmed as differing there. The DKW inflation is markedly slower at the
same nominal PFA, \AblDkwDelay\,s against \AblFullDelay\,s for the exact Beta
levels --- a factor of \AblDkwSpeedup{} --- which quantifies what
\cref{thm:beta} buys conditional on both being valid. And the fixed-start
product martingale (no changepoint prior) misses \AblPointPriorMiss{} of
events: with the change arriving deep into the stream, a wealth process
started at frame one has burned too much to recover, which is precisely the
failure mode the mixture over onsets in \cref{eq:episodic} exists to prevent,
and the structural reason SENTINEL-E outperforms E-SHIFT here: E-SHIFT's
wealth also starts at frame one.

